% answers_short.tex —— 2009 年·全国大纲卷Ⅰ·理科 · 简短答案速查
% 由 tools/make_short_answers.py 自 answers.tex 的 \gksolhead/\gkanswer 提取，
% 并与 2009/2009_数学真题_全国1理_答案要点.md 逐题比对后生成，勿手改（改源 MD 后重跑）。

\begin{gknotice}
\textbf{使用说明}：本节只给最终答案便于快速核对，完整推导见解析版。答案为编写者独立推导，\textbf{非官方标准答案}。
\end{gknotice}

\gkpart{选择题}
\begin{center}\begin{tabular}{ccccccccccccc}
\toprule
题号 & 1 & 2 & 3 & 4 & 5 & 6 & 7 & 8 & 9 & 10 & 11 & 12 \\
\midrule
答案 & A & B & D & C & D & D & D & A & B & C & D & A \\
\bottomrule
\end{tabular}\end{center}

\gkpart{填空题}
\begin{itemize}[leftmargin=2.6em]
  \item[\textbf{13.}] $-240$
  \item[\textbf{14.}] $24$
  \item[\textbf{15.}] $20\pi$
  \item[\textbf{16.}] $-8$
\end{itemize}

\gkpart{解答题}
\begin{itemize}[leftmargin=2.6em]
  \item[\textbf{17.}] $b=4$
  \item[\textbf{18.}] (1) 见解析版；(2) 二面角 $S-AM-B$ 的大小为 $\theta=\arccos\left(-\sqrt{\dfrac23}\right)$
  \item[\textbf{19.}] (1) $\dfrac{81}{125}$；(2) $\Pr(\xi=2)=\dfrac{13}{25}$，$\Pr(\xi=5)=\dfrac{12}{25}$，$\mathbb E\xi=\dfrac{86}{25}$
  \item[\textbf{20.}] (1) $b_n=2-2^{1-n}$；(2) $a_n=2n-\dfrac{n}{2^{n-1}}$，$S_n=n(n+1)-4+\dfrac{n+2}{2^{n-1}}$
  \item[\textbf{21.}] (1) $\dfrac{\sqrt{15}}2<r<4$；(2) $P\left(\dfrac76,0\right)$
  \item[\textbf{22.}] (1) 存在 $x_1\in[-1,0]$、$x_2\in[1,2]$，使 $b=-\dfrac{x_1+x_2}{2}$、$c=x_1x_2$；等价地，$(b,c)$ 是该参数变换下矩形 $[-1,0]\times[1,2]$ 的像；(2) $-10\leqslant f(x_2)\leqslant-\dfrac12$
\end{itemize}
